\sloppy
\section{Discussion}\label{sec:discussion}

\subsection{Why the method works}\label{sec:why}

The posterior $Q$ depends on the covariates only, through $\hat p$ and $\hat Y$; it does not depend on the realised
exposure. Where the estimated exposure propensity is small, $Q$ is close to the imputation $\hat Y$, so sampled
negatives are down-weighted in proportion to the estimated probability that the user would like the item. Where the
estimated propensity is close to one, $Q=\hat Y(1-\hat p)/(1-\hat p\hat Y)$ is close to zero only if $\hat Y$ is also
small. We did not test whether pairs with high estimated exposure propensity are in fact true negatives. Rarely
exposed items are the ones most likely to be false negatives under exposure bias, and this is the motivation for the
weighting. ExpoMF~\citep{liang2016exposure} computes a related posterior for unclicked items
inside matrix factorisation. Here the posterior has a closed form in the nuisances of doubly robust estimation and
enters a pairwise graph loss.

The offset contributes a different signal. It adds a population-level preference estimate, the user and item
effects of an IPS-weighted additive model fitted once on the log. We conjecture that a BPR objective on the observed
positives has difficulty learning such an effect from a biased log, because the exposed items are the ones it sees.
This conjecture was not tested directly. The shuffled-$Q$ control is a check against the explanation that the gain
comes from the variance of the weights alone: a random permutation of $Q$
has the same marginal distribution and removes the gain. On Coat the offset alone reaches essentially the same
accuracy as the combination, so, with unmatched grids and no test, the offset appears to be the larger ingredient there,
while on Yahoo! R3 and KuaiRand the weighting carries most of the gain in the selected runs.

\subsection{When it does not help}\label{sec:when}

KuaiRec is the setting in which the method is weakest. Its log is dense (13.4\% of the matrix is observed) and the
test matrix is fully observed, with 3{,}314 candidates per user. These are the features that differ most from the
other three datasets; we have no direct test of which of them causes the gap. In this setting the pointwise doubly robust
methods (DR-JL and MRDR) and training-free DR propagation are stronger. Pointwise training does not close the gap
either: LightGCN-pn reaches 0.5992--0.5995 and DR-LightGCN-pn 0.5992--0.6027 across the two learning rates, the
latter being the pointwise loss with the same Q weights on negatives. Both are far below DR-JL and MRDR.
On Coat the difference from iALS is not significant across seeds, and iALS has a higher main-run mean than every graph baseline there.

\subsection{Practical guidance}\label{sec:practical}

The following guidance summarises the selections observed in these experiments. It is based on the observed
selections and was not tested as a procedure. For a practitioner with a biased log and a uniform validation sample,
four steps follow from the results. First,
fit the additive IPS imputation and the exposure model on the log, and check that Q is informative with a
shuffled control; if the shuffled model matches the plain model, the weights carry no information in this comparison, and an offset-only model is the natural alternative to test. Second, select the offset weight on validation data and include $\beta=0$ in the grid, because the
validation choice is not always the test-best one (KuaiRand). Third, report the comparison against the strongest
simple baselines, including iALS, because they remain close on some datasets. Fourth, on dense logs with a fully
observed target, compare against training-free DR propagation, which needs no gradient training; on KuaiRec it
reaches 0.6378, against 0.6285 for the trained method.

\subsection{Limitations}\label{sec:limitations}

\begin{enumerate}
\item \textbf{Unequal tuning budgets.} The final method is the union of several jobs. It was tuned with 324
configurations on Coat, against 24 for iALS, 54 for LightGCN and 108 for r-AdjNorm; with 84 on Yahoo! R3 against
72 for r-AdjNorm; with 30 on KuaiRand against 12 for r-AdjNorm; and with 120 on KuaiRec against 6 for DR-JL and
MRDR. Selecting over more configurations can favour the method. The AdvInfoNCE grid is largest on Coat
(216 configurations) and smallest elsewhere (8 or 16), so the comparison with it is unequal in both directions.
\item \textbf{AdvInfoNCE is a simplified re-implementation.} The adversary budget is implicit (the number of
adversarial rounds, with $\eta\in\{4,12\}$), there is no KL-ball projection, and only an embedding-based adversary was
tried. Its numbers are not those of the original paper.
\item \textbf{Small effects.} The KuaiRand gain is $+1.94\times10^{-3}$ in the main run and $+2.31\times10^{-3}$ over
the seeds. The paired differences against the strongest graph baseline range from 0.0013 to 0.0327 nDCG
(Section~\ref{sec:main}), and the confidence-interval half-widths on Coat are between 0.017 and 0.028.
\item \textbf{Main run versus seed runs.} Table~\ref{tab:main} is a single training seed with full selection. The seed
runs retrain only the selected configurations and compare against two baselines per dataset, chosen by main-run test
means. Over these runs the Coat difference from iALS and the KuaiRec differences from DR-JL and MRDR are not significant.
\item \textbf{KuaiRec is lost to the doubly robust baselines.} In the main run the method is below DR-JL and MRDR, and
in all runs it is below training-free DRUP and DR adjacency ($0.6378$).
\item \textbf{One imputation family.} The propensity and the imputation come from one additive model, and on Coat
from one low-rank alternative. Other imputation families are untested.
\item \textbf{Ablations are mostly on Coat.} The offset-only, imputation-only, low-rank and Q-iALS ablations were run
on Coat only. Yahoo! R3 has the popularity and shuffled controls only, and KuaiRand and KuaiRec have only the
decomposition of the selected runs. No significance tests were run for the ablations, except the single KuaiRec
decomposition test reported in Section~\ref{sec:ablations}.
\item \textbf{Coat propensities may leak.} The Coat propensities were estimated by the dataset authors with the uniform
sample~\citep{schnabel2016} that we also use for validation and testing. This is a possible leak that affects Coat only.
The same propensities are used by every propensity-based method on Coat, so the Coat comparison shares this input.
\item \textbf{Exposure semantics.} On Coat and Yahoo! R3, ``exposure'' is self-selection by users rather than platform
exposure. The derivation of $Q$ assumes that exposure is independent of the outcome given the covariates of the
propensity model.
\item \textbf{Overlapping splits.} Per-user $t$-tests over overlapping splits treat the log and the fitted nuisances as
fixed. They can be anti-conservative~\citep{nadeau2003}.
\item \textbf{Scale.} $Q$ and the offset are dense $m\times n$ matrices. The implementation therefore fits the catalogues
used here, the largest being KuaiRec with 7{,}176 users and 10{,}728 items. Computing $Q$ on the fly from the nuisances
would remove this limit; it is not implemented.
\item \textbf{Methods not compared.} StableDR, DICE, CausE, AutoDebias and large language model based recommenders
are not in the comparison. Q-iALS was not run on the larger datasets.
\item \textbf{Validation choice is not always test-best.} On KuaiRand the $\beta=0$ job reaches $0.4623$ on test, against
$0.4613$ for the validated combination.
\end{enumerate}
